\documentclass[11pt, a4paper]{article}

\usepackage[a4paper, top=2.5cm, bottom=2.5cm, left=2cm, right=2cm]{geometry}
\usepackage{amsmath}
\usepackage{amssymb}
\usepackage[colorlinks=true, linkcolor=blue, urlcolor=blue, citecolor=blue]{hyperref}

\title{Theory of Everything - Volume 26 \\ \Large Network Theory \& Graph Dynamics}
\author{Omega Protocol}
\date{\today}

\begin{document}

\maketitle

\section*{Layman's Summary}
Why do the internet, the human brain, and a galaxy of stars all look mathematically similar? Volume 26 proves that the universe naturally organizes itself into a specific kind of network. Reality isn't built on a rigid grid; it's a living, breathing web of connections that follows universal mathematical rules.

\section{Network Topology (Axiom 29)}
We establish that the macroscopic structure of reality can be represented as an expanding "scale-free network." Whether you zoom in on quantum particles or zoom out to cosmic superclusters, you see the exact same web-like pattern. The universe is fundamentally built on connections (edges) between points of information (nodes).

\section{Small World Phenomenon (Theorem)}
Have you heard of "six degrees of separation"? This theorem proves that no matter how gigantic the universe gets, the mathematical "distance" between any two points stays incredibly small. As the network grows, the number of steps required to get from A to B only grows logarithmically. The universe is incredibly tightly knit.

\section{Hub Emergence (Theorem)}
In any network, some nodes become wildly more popular and connected than others (think of Google on the internet, or major airports). We prove that this "preferential attachment" is a mathematical inevitability. Information naturally flows toward nodes that already have information, creating massive "Archive" hubs that anchor the structure of reality.

\end{document}
